% Motivation letter — University of Trento, Bachelor's in Computer Science
% Build: pdflatex motivation-trento-cs.tex (twice)

% Shared LaTeX preamble for all motivation letters.
% Don't compile this file directly --- it's imported by each per-university .tex.
\documentclass[11pt,a4paper]{letter}

\usepackage[a4paper,margin=2.2cm]{geometry}
\usepackage[T1]{fontenc}
\usepackage[utf8]{inputenc}
\usepackage[sfdefault]{roboto}
\usepackage[hidelinks]{hyperref}
\usepackage{xcolor}
\usepackage{microtype}
\usepackage{parskip}

\definecolor{accent}{HTML}{0e7490}

\hypersetup{
  colorlinks=true,
  urlcolor=accent,
}

% Legal name used here — required on university applications.
% Professional name "Ahmad Bagheri" is used on LinkedIn / resume / OSS.
\name{Ahmad Baghereslami}
\address{
  contact@ahmadcodes.com \\
  +37494308785 \\
  ahmadcodes.com \\
  Armenia
}
\signature{Ahmad Baghereslami}

\input{glyphtounicode}
\pdfgentounicode=1


\begin{document}

\begin{letter}{
  International Admissions Office\\
  University of Trento\\
  Via Belenzani 12\\
  38122 Trento, Italy
}

\opening{Dear Admissions Committee,}

I am applying for the Bachelor's in Computer Science (English-taught programme) at the University of Trento for the September 2027 intake. This letter explains who I am, what I have built, and why your programme is the right next step for me.

Last year I built a manufacturer customer portal for Barriertek, a fire-retardant lumber producer in Toronto, replacing phone-tag, whiteboard production scheduling, and per-team Excel sheets across distributors, truckers, production workers, and the office team. End-to-end TypeScript: NestJS API, React + TanStack Query frontend, bi-directional NetSuite integration via SuiteScript RESTlets. PostgreSQL and Linux on a Coolify-managed deployment. I was the sole engineer. The portal is now live with internal users; full external rollout is scheduled for October 2026. I open-sourced the integration client that came out of this work --- \href{https://github.com/Ahmad-b1995/netsuite-restlet-typescript}{github.com/Ahmad-b1995/netsuite-restlet-typescript} --- as my first contribution to the public ecosystem.

Before Barriertek, I led a three-developer frontend team at DexTrading, a cryptocurrency analytics platform, rebuilding the customer application on Next.js with TanStack Query and shipping a NestJS + GraphQL content layer. Earlier I worked at Chargoon, an Iranian ERP vendor, migrating legacy modules to React and .NET alongside a 15-engineer team. Before that I shipped a PWA video streaming product at Hamisheh using Angular and HLS. Across seven years I have accumulated the practical skills a Computer Science Bachelor's programme is meant to teach, but without the structured grounding in algorithms, formal methods, and theoretical computer science that I now want.

My high-school path was Olum-e-Ensani --- the Humanities track in Iran. I never sat for the national Konkur examination because I started programming professionally at eighteen, and the practical path felt more honest than waiting another four years to begin work. That choice shaped seven productive years, but it also means I never received structured mathematics or formal computer-science education. I have studied algorithms and discrete mathematics independently since, but I am clear-eyed about what self-study can and cannot teach. Trento's combination of a rigorous CS core (algorithms, data structures, formal methods) with a strong applied track in software engineering and data systems is exactly the framework I need.

Trento specifically draws me for three reasons. First, the faculty's emphasis on industrial collaboration --- with Microsoft, IBM, and Trentino's regional manufacturing base --- means my professional experience translates naturally into project work rather than sitting on a shelf. Second, the size and international intake match how I work best: smaller cohorts, direct access to professors, focused study rather than the volume of Milan or Turin. Third, Trentino's regional DSU support and Opera Universitaria housing system fit my financial situation as an international student moving with a partner.

During the programme I will continue the open-source work that started with netsuite-restlet-typescript --- building developer tooling for ERP integration that addresses gaps left by the major commercial vendors. My student visa's twenty hours of weekly work I will spend on contract software engagements in my specialty, so my time outside coursework reinforces the trajectory rather than diluting it. After graduation, my goal is to remain in the EU and build a small consulting practice in ERP integration for traditional industrial businesses. Italy and Germany both have manufacturing sectors that need exactly this kind of technical bridging, and Trento --- positioned between the two --- is an unusually good vantage point for that work.

I would be grateful for the opportunity to formalize seven years of self-directed practice under your faculty.

\closing{Sincerely,}

\end{letter}

\end{document}
